\chapter*{Abstract}
\addcontentsline{toc}{chapter}{Abstract}

This book presents a comparative approach to PID controller tuning in isolated power
systems, using a hybrid hydro-solar microgrid in Limbani, Puno, as the reference case.
The study starts from the well-known limitations of classical tuning rules when they are
applied to plants affected by nonlinear dynamics, water hammer, low inertia, and severe
disturbances caused by load changes or solar intermittency.

The tuning task is formulated as an optimization problem in which the controller gains
$K_p$, $K_i$, and $K_d$ are searched within constrained domains according to dynamic
performance criteria such as overshoot, settling time, and integral error indices, with
special emphasis on the ITAE criterion. On that basis, the book reviews the fundamentals
of metaheuristic optimization and discusses several search strategies, highlighting Particle
Swarm Optimization (PSO) because of its balance between implementation simplicity,
computational effort, and convergence quality.

The comparative discussion shows that optimized tuning can improve frequency recovery,
reduce sustained oscillations, and decrease the mechanical stress imposed on hydraulic
actuators. The text also outlines how this framework can be extended toward additional
swarm-based methods, hybrid optimizers, and more realistic validation environments.

\bigskip
\noindent\textbf{Keywords:} metaheuristics, PID tuning, isolated microgrids, optimization, PSO.
